RailFlux Route Assignment Procedure - Comprehensive Technical Analysis
Request Initialization and Entry Framework
Route Request Entry Point: RouteAssignmentService.requestRoute method establishes primary interface accepting sourceSignalId, destSignalId, direction specification, requestedBy operator identification, trainData QVariantMap containing train characteristics, and priority classification (LOW/NORMAL/HIGH/EMERGENCY). System-wide request counter increments immediately establishing comprehensive statistical tracking throughout route lifecycle management.
System Operational Validation: Entry point verification includes operational status checking through m_isOperational boolean flag, service health validation through areServicesHealthy method encompassing DatabaseManager connectivity, GraphService loading status, ResourceLockService operational state, OverlapService functional status, TelemetryService availability, and VitalRouteController operational verification ensuring complete service composition functionality before proceeding with route processing.
Request Capacity Management: canAcceptNewRequests method evaluates system capacity through operational status verification, request queue size comparison against MAX_QUEUE_SIZE threshold (50 requests maximum), and emergency mode exclusion preventing new route acceptance during critical system states. Capacity violation triggers systemOverloaded signal emission with queue size and maximum concurrent route parameters enabling UI feedback and operational awareness.
RouteRequest Object Construction: Request object assembly includes QUuid generation through createUuid method for unique identification, timestamp assignment via currentDateTime for chronological tracking, priority assignment for queue management, operator identification for audit trails, train data preservation for overlap calculations, and reason assignment establishing request context for operational analysis and audit compliance.
Priority-Based Queue Management Architecture
Multi-Criteria Priority Scoring: calculateRequestPriority method implements hierarchical scoring algorithm beginning with base score of 100 points, emergency priority adding 1000 points ensuring absolute precedence, high priority contributing 500 points for operational importance, normal priority maintaining base scoring, low priority subtracting 100 points for deferrable operations, and age factor incrementally adding one point per ten seconds ensuring temporal fairness preventing request starvation.
Dynamic Queue Prioritization: prioritizeQueue method converts queue to list structure enabling sophisticated sorting operations, applies multi-criteria comparison function utilizing priority score primary sorting with higher values preceding lower values, implements timestamp-based secondary sorting with earlier requests taking precedence among equal priority requests, and rebuilds queue maintaining prioritized order ensuring optimal processing sequence throughout system operation.
Overload Detection and Management: addToQueue method monitors queue size against OVERLOAD_THRESHOLD (20 requests) triggering automatic system protection, emits systemOverloaded signal providing queue size and maximum concurrent route information for operator awareness, records safety event through TelemetryService with WARNING severity classification documenting system stress conditions, and maintains operational statistics enabling performance analysis and capacity planning optimization.
Queue Processing Control: shouldProcessRequest method evaluates concurrent route limitations through activeRoutes count comparison against maximum permitted routes (10 normal operation, 5 degraded mode), system degradation status assessment modifying operational parameters for conservative resource management, and operational status verification ensuring safe processing conditions before dequeueing next route request for pipeline execution.
Five-Stage Processing Pipeline Architecture
Stage One - Comprehensive Validation: validateRequest method delegates to VitalRouteController.validateRouteRequest implementing safety-critical validation through signal existence verification, operational status confirmation, interlocking rule compliance checking, resource availability assessment, signal type compatibility validation, and operator authorization verification. Validation failure immediately terminates processing with descriptive error messaging and performance metrics recording enabling operational analysis and troubleshooting capabilities.
Stage Two - Intelligent Pathfinding: performPathfinding method implements signal-to-circuit resolution through resolveSignalToCircuit extracting succeededByCircuitId from source signals and precededByCircuitId from destination signals, invokes GraphService.findRoute with start circuit, goal circuit, direction constraints, current point machine states, and configurable timeout parameters (500ms maximum), processes A-star algorithm results including path discovery, cost calculations, node exploration statistics, and processing time metrics establishing comprehensive pathfinding intelligence.
Stage Three - Safety Overlap Computation: calculateOverlap method delegates to OverlapService.calculateOverlap processing source signal, destination signal, direction parameters, and train data characteristics enabling dynamic overlap determination. Overlap calculation encompasses fixed overlap retrieval from signal definitions, dynamic overlap adjustment based on train speed and braking characteristics, safety margin application for conservative operation, hold time determination for automatic release scheduling, and method classification (FIXED/DYNAMIC/SAFETY_MARGIN) supporting operational analysis and safety compliance.
Stage Four - Intelligent Route Establishment: System transitions from traditional resource reservation to VitalRouteController.establishRouteWithIntelligentAspects method implementing AspectPropagationService integration for sophisticated signal aspect coordination, point machine position optimization based on pathfinding results, control dependency analysis ensuring signal interlocking compliance, resource conflict resolution through intelligent alternative discovery, and coordinated infrastructure execution maintaining operational safety and efficiency.
Stage Five - Database Persistence and Finalization: finalizeRoute method implements route assignment persistence through database insertion, overlap resource reservation via OverlapService.reserveOverlap when safety margins require protection, route state initialization establishing lifecycle tracking, performance metrics recording for operational analysis, and success confirmation enabling route activation and operator notification through comprehensive signal emission.
Intelligent Aspect Propagation Integration
AspectPropagationService Integration: VitalRouteController.establishRouteWithIntelligentAspects method leverages control graph construction identifying signal dependencies, performs graph pruning focused on destination-relevant control paths, creates dependency ordering through topological sorting ensuring proper signal progression sequences, applies forward aspect propagation with intelligent selection algorithms, and validates destination constraints ensuring stopping points receive appropriate RED aspects or proceed indications based on track clearance conditions.
Control Dependency Processing: Intelligent propagation evaluates control relationships through InterlockingRuleEngine.validateInterlockedSignalAspectChange processing current aspects, requested aspects, point machine conditions, and control mode specifications (AND requiring all controllers to permit, OR requiring any controller to permit). System respects signal independence for autonomous operation while enforcing dependency constraints for controlled signals ensuring comprehensive safety compliance throughout route establishment procedures.
Point Machine Coordination: Aspect propagation incorporates point machine position requirements derived from pathfinding analysis, validates required position changes through isPointMachineSettable verification including operational status, lock conditions, protecting signal verification, and mechanical availability assessment. Coordinated execution ensures simultaneous signal aspect changes and point machine movements preventing intermediate unsafe configurations during route establishment.
Coordinated Execution Framework: establishRouteWithIntelligentAspects executes systematic infrastructure changes through signal aspect updates using intelligent selection algorithms, point machine position changes based on pathfinding requirements, resource lock acquisition through ResourceLockService preventing conflicts, database state synchronization maintaining consistency, and comprehensive error handling with rollback capabilities ensuring atomic operations throughout route establishment procedures.
Resource Conflict Detection and Management
Multi-Type Resource Locking: ResourceLockService implements comprehensive locking mechanisms supporting EXCLUSIVE locks preventing any concurrent access, SHARED locks permitting compatible operations, OVERLAP locks restricting safety margin areas during route overlap periods, and EMERGENCY locks overriding standard conflict resolution enabling critical operations during system emergencies. Lock acquisition includes timeout management, conflict detection algorithms, deadlock prevention strategies, and comprehensive audit logging for regulatory compliance.
Track Circuit Reservation System: Circuit locking encompasses occupancy verification through database queries confirming clear status, reservation conflict detection identifying existing route allocations, resource availability assessment ensuring mechanical and electronic system readiness, and lock duration management with automatic expiration preventing indefinite resource retention. Database integration includes resource_locks table management with route identification, resource type classification, acquisition timestamps, operator tracking, and release reason documentation.
Point Machine Lock Coordination: Point machine locking includes paired machine coordination ensuring synchronized operations, protecting signal verification requiring RED aspects before movement authorization, mechanical status assessment confirming operational readiness, and safety interlock validation preventing unauthorized operations during train movements. Lock management encompasses position consistency verification, transition state monitoring, failure condition handling, and maintenance lock recognition preventing operational interference.
Conflict Resolution Algorithms: Resource conflict detection implements comprehensive analysis including active route identification, resource overlap assessment, lock type compatibility evaluation, temporal conflict analysis, and alternative resource discovery. Resolution strategies include automatic conflict avoidance through intelligent resource selection, operator notification for manual intervention, and emergency override capabilities supporting critical operational requirements while maintaining safety compliance.
Dynamic Safety Overlap Management
Overlap Calculation Intelligence: OverlapService.calculateOverlap implements multiple calculation methods including fixed overlap retrieval from signal definitions providing baseline safety margins, dynamic overlap computation incorporating train characteristics including speed, length, braking rates, and track conditions, and safety margin overlap applying conservative multipliers for enhanced protection during adverse conditions. Calculation results include overlap circuit identification, hold time determination, release trigger specification, and method classification supporting operational analysis.
Train-Specific Overlap Adjustment: Dynamic overlap calculation processes train data including speed conversion from kilometers per hour to meters per second, braking distance computation using physics-based calculations (v²/2a), safety margin addition incorporating train length and buffer distances, and hold time optimization balancing safety requirements with operational efficiency. Speed factor normalization enables consistent overlap scaling across varying operational conditions while maintaining safety standards.
Overlap Resource Management: reserveOverlap method implements comprehensive overlap resource control including circuit availability verification, conflict detection with existing reservations, lock acquisition through ResourceLockService integration, automatic release timer configuration, and trigger condition monitoring. Resource management includes database persistence through overlap state tracking, performance metrics collection, and operational statistics maintenance supporting system optimization and regulatory compliance.
Automatic Release Mechanisms: Overlap release triggers include track circuit clearance detection indicating train passage beyond overlap boundaries, timer-based automatic release providing time-bounded safety margins, operator manual release enabling operational flexibility, and emergency release supporting critical system recovery. Release validation includes safety condition verification, resource availability confirmation, and audit logging ensuring comprehensive operational tracking and compliance documentation.
Performance Monitoring and Telemetry Integration
Comprehensive Metrics Collection: TelemetryService.recordPerformanceMetric captures detailed processing statistics including stage-specific timing (validation, pathfinding, overlap calculation, resource reservation, finalization), success/failure rates with categorized failure analysis, resource utilization tracking across track circuits and point machines, and system performance indicators including queue depths, processing latencies, and throughput metrics supporting operational optimization and capacity planning.
Safety Event Documentation: TelemetryService.recordSafetyEvent implements comprehensive safety monitoring including event type classification (violations, warnings, critical failures), severity level assignment (INFO, WARNING, ERROR, CRITICAL), entity identification specifying affected infrastructure, contextual information providing detailed failure analysis, and operator correlation enabling accountability tracking and training identification for enhanced system safety.
Real-Time Performance Analysis: System implements continuous performance monitoring including stage performance tracking maintaining rolling averages for validation, pathfinding, overlap calculation, and resource reservation phases, threshold violation detection triggering automatic alerts when processing times exceed WARNING_PROCESSING_TIME_MS (2000ms) or queue sizes exceed capacity thresholds, and statistical analysis providing operational insights supporting system optimization and preventive maintenance scheduling.
Operational Statistics Compilation: Performance analysis encompasses total request counting, successful route establishment rates, failure categorization with root cause analysis, emergency release frequency indicating system stress conditions, timeout occurrence tracking revealing performance bottlenecks, and success rate calculation providing overall system effectiveness metrics. Statistics support operational decision-making, maintenance scheduling, and system capacity optimization throughout continuous railway operations.
Database Integration and State Management
Route State Machine Implementation: Route lifecycle management encompasses state transitions from REQUESTED through VALIDATING, RESERVED, ACTIVE, PARTIALLY_RELEASED, to final RELEASED state, with alternative paths including FAILED, EMERGENCY_RELEASED, and DEGRADED states supporting comprehensive operational scenarios. Database persistence through route_assignments table maintains complete lifecycle tracking with timestamp precision, operator identification, and performance metrics preservation.
Atomic Database Operations: persistRouteAssignment method implements comprehensive database integration including route parameter validation, signal existence verification, circuit availability confirmation, transaction management ensuring atomic operations, constraint validation preventing data inconsistency, and audit logging establishing complete operational traceability. Database functions include sophisticated parameter validation, error handling with descriptive messaging, and performance optimization through prepared statement utilization.
Event Sourcing Architecture: Route lifecycle tracking through route_events table captures comprehensive operational history including event type classification, timestamp precision, operator identification, contextual data preservation, performance metrics, and safety significance flags. Event sourcing enables complete operational reconstruction, forensic analysis capabilities, regulatory compliance documentation, and system behavior analysis supporting continuous improvement initiatives.
Real-Time Notification System: Database triggers automatically generate PostgreSQL NOTIFY events through pg_notify function calls enabling immediate application-level updates, QML property binding updates, and operator interface synchronization. Notification processing includes payload parsing, entity identification, change classification, and appropriate signal emission maintaining system consistency throughout distributed operational components.
Emergency Handling and System Recovery
Emergency Mode Activation: activateEmergencyMode triggers comprehensive system protection including degraded mode engagement reducing concurrent route limits, conservative timeout application extending processing intervals, non-essential request clearing maintaining only EMERGENCY and HIGH priority operations, and safety event logging documenting emergency conditions with complete contextual information supporting post-incident analysis and system recovery procedures.
Emergency Route Release: processEmergencyReleaseAll implements systematic route termination including active route identification, resource lock release for all affected infrastructure, overlap termination regardless of hold timers, database state synchronization ensuring consistency, and comprehensive audit logging documenting emergency actions with operator identification, timing precision, and contextual reasoning supporting regulatory compliance and operational analysis.
System Health Monitoring: checkSystemHealth continuously evaluates service availability including DatabaseManager connectivity verification, GraphService operational status, ResourceLockService functional assessment, OverlapService availability confirmation, TelemetryService operational verification, and VitalRouteController safety system status. Health degradation triggers automatic protection measures including service isolation, degraded mode activation, and operator notification ensuring system stability throughout adverse conditions.
Recovery Procedures: System recovery encompasses service restoration through health verification, configuration reloading, queue processing resumption, performance monitoring restoration, and operational status communication. Recovery validation includes comprehensive testing, performance verification, safety system confirmation, and gradual capacity restoration ensuring safe system operation throughout recovery processes while maintaining operational continuity and safety compliance.